\documentclass[11pt,a4paper]{ctexart}
\usepackage[margin=2.2cm]{geometry}
\usepackage{amsmath,amssymb,booktabs,float,xcolor,hyperref,enumitem,xurl}
\hypersetup{colorlinks=true,linkcolor=blue!50!black,urlcolor=blue!50!black}
\setlist{itemsep=2pt,topsep=3pt}
\title{第 01 题的 C/MPI 并行实现（bt-v + 混合精度迭代修正）}
\author{温州肯恩大学 \quad Nonlinear Poisson 项目}
\date{2026 年 11 月 4 日}

\begin{document}
\maketitle

\begin{abstract}
把 H41 的方案（电子消元块三角预条件 bt-v + 双双精度迭代修正）写成了 C + MPI + hypre 程序，直接嵌入官方示例框架（\texttt{problems/problem01/solver.c}）。
在真实矩阵上：4 个进程端到端 7.5--8.4\,s（框架的 5 次计时平均，包含分析、准备和求解），单进程 21.3\,s；Python 版本单线程为 34.3\,s。
真实残差 $\|b-Ax\|/\|b\|=8.5\times10^{-11}$，用独立的 Python 双双精度程序复核结果完全一致；内层迭代数（20+9+10）与 Python 版本（19+9+10）一致。
总内存 4.2--5.2\,GB，低于 9\,GB 限制。笔记本当时负载很高（负载平均最高约 196），超过 4 个进程后反而变慢，真正的并行扩展必须在比赛机器上测。
float128 输出路径在 macOS 上无法编译测试（没有 \texttt{\_\_float128}），需要在 Linux/GCC 上验证。
\end{abstract}

\section{实现}
代码在 \texttt{FD\_LGF/\allowbreak solver\_challenge/\allowbreak c\_impl/}（官方框架的副本），我们只改了 \texttt{problems/problem01/} 下的文件：
\begin{itemize}
\item \texttt{solver.h / solver.c}（约 590 行）：框架的 \texttt{analyse}、\texttt{preprocess}、\texttt{iterative\_solver} 三个接口。
\item \texttt{CMakeLists.txt}：链接 hypre（\texttt{-DHYPRE\_DIR=...}），并加 \texttt{-ffp-contract=off}（双双精度需要乘法不被融合）。
\item \texttt{README\_problem01.md}：比赛机器上的编译和运行方法；本地用 \texttt{build\_local.sh}（conda 的 OpenMPI + hypre 3.2）。
\end{itemize}
\paragraph{分析阶段（analyse）。}框架按行均分矩阵，但三个温度的同一单元落在不同进程上。我们先用 \texttt{MPI\_Alltoallv} 把每行发给“该单元的所有者”，
使每个进程拥有一段连续单元的全部三个温度；然后建立本地 CSR（本地列 + 邻居列）和邻居交换表。
\paragraph{准备阶段（preprocess）。}取出对角耦合 $a_{12},d_{22},a_{23},a_{32}$；用 hypre IJ 接口建立辐射块 $A_{11}$ 和离子 Schur 补 $S_3$（对角元按无抵消公式），
各做一次 BoomerAMG 设置（HMIS、ext+i、P\_max 4、$\theta=0.5$、一层激进粗化、对称混合 Gauss--Seidel、最粗层直接解，与 Python/PETSc 版本相同）。
\paragraph{求解阶段（iterative\_solver）。}外层迭代修正：$x$ 以双双精度（约 32 位有效数字）保存，残差 $b-Ax$ 用双双精度逐行计算（乘法误差用 \texttt{fma} 精确求出）；
内层为自写的右预条件 GMRES(50)（修正 Gram--Schmidt，$10^{-6}$），预条件为 bt-v（两次标量 V 循环 + 对角运算）。为省内存，只存 Krylov 基 $V$，每次重启结束时再做一次预条件得到修正量。
最后把解送回框架的行划分；用 \texttt{SC26\_SOLUTION\_FP128} 编译时，解以 \texttt{\_\_float128} 的形式交给框架，官方的 float128 检查就能看到真实的 $10^{-10}$ 级残差。
\paragraph{参数。}环境变量 \texttt{SC26\_P01\_TOL} 设目标（默认 $10^{-8}$；第一档用 $10^{-4}$），\texttt{SC26\_P01\_INNER} 设内层容差，另有输出每步信息和导出解向量的诊断开关。

\section{正确性}
\begin{itemize}
\item 4 进程运行，每步残差（双双精度）：$1.0\to1.35\times10^{-2}\to6.4\times10^{-6}\to8.5\times10^{-11}$，内层 GMRES 20、9、10 步。Python 版本为 19、9、10 步，残差 $8.0\times10^{-11}$。
\item 导出的双双精度解用独立的 Python 程序（\texttt{check\_dd\_solution.py}，基于此前与有理数精确计算核对过的 \texttt{ir\_fp128.DDResidual}）复核：$8.513\times10^{-11}$，与 C 程序一致。
\item 把同一个解只取高位（即普通双精度解）：残差 $6.4\times10^{-3}$；框架的双精度检查显示 $8.4\times10^{-3}$。这再次说明必须用 float128 输出。
\end{itemize}

\section{时间与内存}
macOS 笔记本（6 个性能核 + 2 个能效核），OpenMPI，每个进程单线程。机器上同时有其他程序在运行（负载平均 5--196），计时只能作参考。
\begin{table}[H]
\centering
\small
\begin{tabular}{rrrrr}
\toprule
进程数 & 目标 $10^{-8}$（s） & 目标 $10^{-4}$（s） & 总内存（GB） & 当时负载 \\
\midrule
1 & 21.3 & 16.5 & 4.2--5.0 & 5 / 110 \\
2 & 12.8 & 10.1 & 4.2--5.0 & 19 / 11 \\
4 & 7.5，7.5，8.0，8.4 & 7.5 & 4.2--4.9 & 4--196 \\
6 & 9.2，10.7 & 9.2 & 4.2--5.0 & 8--73 \\
8 & 11.9 & 11.9 & 5.1--5.2 & 9 / 42 \\
\bottomrule
\end{tabular}
\caption{框架报告的端到端时间（5 次预热后 5 次平均，含分析、准备、求解）。4 进程的详细分解：分析 + 准备约 2.8\,s，求解约 5.2\,s（3 步修正）。
内存为各进程峰值之和（框架在 macOS 上把字节当作 KB 打印，表中已换算）。}
\end{table}
\begin{table}[H]
\centering
\small
\begin{tabular}{p{7.2cm}rp{6cm}}
\toprule
版本 & 时间 & 说明 \\
\midrule
Python/PETSc：bt-v + 迭代修正（H42 批次） & 34.3\,s & 单线程，真实残差 $8.0\times10^{-11}$ \\
\textbf{C/MPI：同一算法，1 进程} & \textbf{21.3\,s} & 同一代码（此次未导出解单独复核） \\
\textbf{C/MPI：4 进程} & \textbf{7.5--8.4\,s} & 真实残差 $8.5\times10^{-11}$ \\
参考：Python/PETSc 整体 BoomerAMG + GMRES（只到自报 $10^{-8}$） & 18--25\,s & 单线程，真实残差停在 $1.3\times10^{-2}$，达不到官方判据 \\
\bottomrule
\end{tabular}
\end{table}
观察：
\begin{itemize}
\item 1 $\to$ 2 $\to$ 4 进程加速 1.7 和 2.8 倍；6、8 进程反而变慢，原因是能效核较慢、机器被其他程序占用，并且每个进程的 MPI 通信相对计算变大。笔记本上无法判断真实的扩展性。
\item $10^{-4}$ 只需 2 步修正，理论上应比 $10^{-8}$ 少约 1.2\,s；在 4 进程上看不出差别，说明计时噪声与这个差值同量级。
\end{itemize}

\section{还没有验证的部分}
\begin{enumerate}
\item \textbf{float128 路径：}\texttt{SC26\_SOLUTION\_FP128} 编译和官方 float128 检查需要 Linux + GCC（libquadmath）。代码已按框架的类型写好，但没有实际编译运行过。
\item \textbf{规则：}float128 解向量、以及两档判据分别用 \texttt{SC26\_P01\_TOL} 运行，是否被允许，需要问组委会。如果两档必须用同一次运行，就只能都按 $10^{-8}$ 求解。
\item \textbf{比赛机器上的扩展性和内存：}96 核 Xeon（或 608 核 LX2）上需要重新测；进程多时，每个进程的 MPI/hypre 固定开销会增加总内存，应确认仍低于 9\,GB。
\item \textbf{GPU（DCU）版本：}未做。
\end{enumerate}

\section{下一步建议}
\begin{enumerate}
\item 在任一台 Linux/GCC 机器上编译 FP128 版本，跑一次官方 float128 检查，确认得分路径打通。
\item 上比赛机器测 24/48/96 进程的时间与内存；视结果调整进程数和 GMRES 重启长度（内存紧时可减到 30）。
\item 性能优化：把修正量的预条件与 Krylov 基的线性组合合并、用一次归约的 CGS2 代替修正 Gram--Schmidt（减少全局通信）、内层改用单精度 AMG。
\item 算法上仍可继续：同时改进两个标量块（H42 的结论），目标是外层 10 步以内且每步开销不超过 bt-v 的 1.5 倍。
\end{enumerate}
\end{document}
